% NAL 2339, batch 3 — Gallica views 41–60.
% Pass: Opus 5.5 (claude-opus-5-5), 2026-09-23 — first pass, unchecked against the views by a human.
%
% The views. Colour portrait scans, one page per view, about 3830 × 5450
% (v. 41–54), larger for the end leaves (to 4300 × 5900). Read from the local
% mirror only (archives/nal-2339), as four full-width bands at 2400 px and
% tight crops for doubtful words; nothing was fetched from Gallica.
%
% What the twenty views hold. One hand throughout v. 41–51, a posed
% seventeenth-century cursive with great flourished initials, brown ink,
% line-ends filled with strokes: a fair copy, whether the author's or a
% copyist's the leaves do not say. Attribution: the catalogue's title names
% Fermat; so do the leaves, in two headings (v. 41 « A Domino de Fermat ad
% Dominum De Carcaui », v. 49 « Lettre de Mons.r De fermes A Monsieur De
% Carcaui »). Neither heading says the hand is Fermat's, there is no
% signature, and nothing here is attributed beyond what these headings carry.
%
%   v. 41–48  A Latin piece, titled on v. 41 « [Nouus] Secundarum et
%             Vlterioris ordinis Radicum In Analyticis Vsus A Domino de
%             Fermat ad Dominum De Carcaui Die 20a. Aprilis Anno 1650
%             Missus ». It begins at the head of v. 41 (so v. 41 does not
%             continue a text from v. 40, as far as this leaf shows). The
%             method: with two unknowns A (first root) and E (second root),
%             collect the terms affected by E on one side of each of two
%             equations, set the two equations as a proportion (« duplicatæ
%             æqualitatis analogia »), cross-multiply, divide by E, and
%             iterate until E appears only « sub Latere », when it is given
%             by a single division (« Applicatio »); worked example
%             A³ + E³ = Z solid, BA + E² + DE = N², carried to a fraction
%             for E (v. 41–43); extension to third, fourth … roots (v. 43).
%             v. 44–46 « Appendix Ad superiorem Methodum »: the method
%             removes all irrationals (« asymmetriæ ») — a sum of cube,
%             square, fourth roots equal to a line; Viète's analyst is named
%             (« Analysta Vietæus ») as defeated by it; worked example
%             Lat. Cub. (ZA² [− A³]) + L. Cub. (A³ + B²A) = D, putting
%             Lat. Cub. (A³ + B²A) = E and cubing (v. 44–45); then
%             determinate problems (one unknown position), local problems
%             (two: a line), loci ad superficiem (three), and « abundant »
%             problems with more data than needed (v. 46–47), with an
%             « Exemplum » (a segment ABC, rectangle and difference of
%             squares given) and the species example A³ + B²A = Z²D,
%             GsA − A⁴ = B²N pl. (v. 47); v. 48: cutting a cone with given
%             vertex and elliptic base in a circle, by six points instead of
%             five; diameters and axes of curves and loci ad superficiem;
%             tangents mentioned last. Ends on v. 48 with a closing flourish.
%   v. 49–51  A French letter, headed « Lettre de Mons.r De fermes / A
%             Monsieur De Carcaui De Castres / Le 20e auril 1650. »,
%             beginning « Monsieur Je n'ay point voulu differer a vous
%             Enuoyer ma methode generale pour le desbrouilement des
%             Asymmetries »: the writer sends « mon original », asks for it
%             back or a copy, claims the discovery; proposes a tangent
%             problem to a curve given by a sum of cube, fourth, square and
%             fifth roots (the formula in Latin, v. 49); says the method of
%             « Mons.r Des Cartes » is more encumbered; « ma methode de
%             Tangentibus » gives the tangent by « la seule application »;
%             Descartes « proposoit comme vne difficulté Insurmontable » the
%             tangent to a curve whose distances to four given points have a
%             given sum, « Ainsy que Je puis faire voir par vne de ses
%             Lettres »; « En voila trop pour vne seconde Lettre »; « mon
%             escrit Latin ». Ends on v. 51 with « &c. », no closing formula,
%             no signature. The letter and the Latin piece bear the same date
%             (20 April 1650); that the letter covers the piece is an
%             inference from its wording.
%   v. 52     Verso of v. 51: blank, the text of the recto showing through.
%   v. 53–56  Two blank leaves (rectos numbered 23 and 24 in red, see
%             below) and their blank versos; v. 56 shows only the offset of
%             a stamp and a few specks.
%   v. 57–59  Blank modern endleaves.
%   v. 60     Marbled endpaper.
%   Views 52–60 are not transcribed. Nothing runs on past v. 51.
%   v. 61–66 (batch 4): endpaper, back board, spine lettered MÉMOIRES DE
%   FERMAT, and the edges; not transcribed, and batch 4 has no file.
%
% Numbers on the leaves. Two series stand one above the other at the top
% right of every recto, and agree: a thin grey pencil figure (17, 18, 19,
% 20, 21., 22 on v. 41, 43, 45, 47, 49, 51) and below it a larger red
% crayon figure with the same number (the 19 and 22 underlined in red).
% The red series goes on over the blank leaves (23 on v. 53, 24 on v. 55),
% where no grey figure was seen. Versos carry none; each shows its recto's
% number reversed through the paper. The grey pencil series is one number
% per leaf, thin, top right of rectos only, in no writer's hand: it is taken
% as the library's foliation and written \folio{17r}–\folio{22r}. The red
% series duplicates it and counts the blank leaves too; it looks like a
% later library count and is noted, not used. The coordinator should check
% that batches 1–2 read the same pencil series (1–16) before these \folio{}
% stand; this batch saw no certificate. No pagination of the writer's was
% seen. Round stamps « BIBLIOTHÈQUE NATIONALE R.F. » on v. 41 and v. 51.
%
% Dating. The catalogue gives none; \dating{} is left empty and no date is
% inferred for the volume. The only dates are on the leaves: « Die 20a.
% Aprilis Anno 1650 » (v. 41) and « Le 20e auril 1650 » (v. 49), given where
% they stand.
%
% Notation and spelling. The species notation of the leaves is kept: A, E
% the unknowns, consonants the givens, dimension by an abbreviation in
% superscript (Z^s, Z^so = Z solido; N q^dto = N quadrato; N pl. = N plano;
% B pp^l = B plano-plano; G^s A = G solido in A), powers as numeral exponents
% (A³, E², N⁴), products by juxtaposition (BA, DE, BZ^sA), roots as « Lat.
% Cub. (…) », « Lat. quad. quad. (…) », « Lat quad. Cub. (…) ». Equality is
% an abbreviation, never a sign: « Æle » (æquale), « Æ. », « Æ^tur »
% (æquatur), « Æquabit^ur », kept as written in text between formulas. The
% proportion of v. 42 separates its terms with « / » and « // », kept.
% Latin and French as written: u/v as on the leaf (« Nouus », « deuoluatur »,
% « Vt », « vne », « auril »), the ligature æ/Æ as written (the hand makes
% it capital-sized in « Æquari », « Æquatur » and lower-case elsewhere),
% capitals inside sentences as written (« In », « Ipse », « Vnicam »), the
% writer's slips (« æqualitis », « Redductio », « asendat », « Extration »)
% kept. Abbreviations by a bar over the word are given with a tilde over the
% vowel (« raõe » ratione, « determinaõne », « operaõn », « applicaõne »,
% « rectangulũ ») and expanded in notes. The long s is set s. No glyph
% lacks a LaTeX equivalent. A slip in the calculation (v. 42, N²E² for N²E³)
% is left and noted.
%
% Continuities. v. 41 opens a new titled piece; whether v. 40 ends another
% cannot be told from this batch. v. 51 ends the letter; v. 52–60 are blank
% or binding, so nothing runs on into batch 4.

% Coordinator's reconciliation (2026-09-23), after all three batches:
%   The red crayon series, top right of every recto, blanks included, runs
%   1–6 (v. 9–19), 7–16 (v. 21–39), 17–24 (v. 41–55) without a gap, and ends
%   on the two blank leaves 23 and 24, exactly as the certificate of v. 7
%   has it (« Volume de 24 Feuillets / Les Feuillets 23 & 24 sont blancs »).
%   It is the library's foliation, and not ink: \folio{} is written from it
%   on every transcribed recto, 1r–22r. The grey pencil series beside it is
%   one higher on v. 9–19, irregular on v. 21–39, and agrees from v. 41; it is
%   recorded in notes and not used. Batches 1 and 2 first wrote no \folio{};
%   the coordinator added them. The ink pagination of the loci treatise
%   (1–18) and the foot numbers 43–45 of the Appendix are the copies' own.
%   v. 61–66 (batch 4) are binding, so batch 4 has no file.
%   Fermat is named on the leaves in three headings only: « A Domino de
%   Fermat ad Dominum De Carcaui » (v. 41), « Lettre de Mons.r De fermes A
%   Monsieur De Carcaui » (v. 49), and « Mr. fermat » over v. 37 in a paler,
%   perhaps later hand; the certificate's « Mémoires de Fermat » (v. 7) is
%   the library's. No hand is attributed.

\documentclass[11pt,a4paper]{article}
% The preamble shared by every transcription of the mathematicians' manuscripts in Gallica.
%
% It rests on one idea: **uncertainty compounds, it does not resolve.** A
% transcription that smooths over a doubtful word has destroyed the only thing
% it was made for — a reader must be able to tell, at every point, what was on
% the page and what was supplied. The apparatus macros below are therefore the
% critical apparatus, not decoration.
%
% The same commands are recognised by `scripts/render.mjs`, which produces the
% site's reading view, and by `scripts/tei.mjs`. A macro added here must be
% added there too, or rendering fails loudly — which is the wanted behaviour.
%
% Comments here are English, like the rest of the repository. The documents
% this preamble sets are French, and that is deliberate: see the skills.

\ProvidesPackage{germain}[2026/09/15 Transcriptions of mathematicians' manuscripts in Gallica]

\usepackage{amsmath,amssymb,amsthm}
\usepackage{mathrsfs}
% Kept from the parent preamble so the renderer's diagram support stays
% available; these manuscripts draw no commutative diagrams, and nothing here uses it.
\usepackage{tikz-cd}
\usepackage[english,french]{babel}
\usepackage{fontspec}

% --- Greek ------------------------------------------------------------------
% The text face stays fontspec's default, Latin Modern, which has no Greek:
% XeTeX drops every glyph it lacks with a line in the log and nothing on the
% page, and the Greek words the manuscripts quote vanished from the PDFs. So
% the Greek and Greek Extended blocks get a XeTeX character class of their own,
% and entering or leaving a run of it switches the family to CMU Serif — the
% same Computer Modern design, with polytonic Greek — and back. Only the
% family changes: a Greek word in \textit or \textbf keeps its shape and series.
% The transitions to and from every other class carry the tokens that class
% has towards an ordinary letter, so French spacing before « ; » or inside
% « » still holds after a Greek word. No grouping, so no brace or \endgroup
% can come out unbalanced; maths is left alone, since XeTeX inserts nothing
% there. Kept identical in `exercices.sty`.
\makeatletter
\newfontfamily\germainGreekfont{cmun}[
  Extension=.otf, NFSSFamily=germaingreek,
  UprightFont=*rm, ItalicFont=*ti, SlantedFont=*sl,
  BoldFont=*bx, BoldItalicFont=*bi, BoldSlantedFont=*bl]
\newXeTeXintercharclass\germain@greek
\def\germain@greekrange#1#2{%
  \@tempcnta=#1\relax
  \loop
    \XeTeXcharclass\@tempcnta=\germain@greek
    \ifnum\@tempcnta<#2\advance\@tempcnta\@ne\repeat}
\germain@greekrange{"0370}{"03FF}
\germain@greekrange{"1F00}{"1FFF}
\def\germain@greekprev{\rmdefault}
\def\germain@greekin{%
  \edef\germain@greekprev{\f@family}\fontfamily{germaingreek}\selectfont}
\def\germain@greekout{\fontfamily{\germain@greekprev}\selectfont}
\def\germain@greekjoin{%
  \XeTeXinterchartokenstate=\@ne
  \@tempcnta=\z@
  \loop
    \ifnum\@tempcnta=\germain@greek\else
      \XeTeXinterchartoks\germain@greek\@tempcnta=\expandafter{%
        \expandafter\germain@greekout\the\XeTeXinterchartoks\z@\@tempcnta}%
      \XeTeXinterchartoks\@tempcnta\germain@greek=\expandafter{%
        \the\XeTeXinterchartoks\@tempcnta\z@\germain@greekin}%
    \fi
    \ifnum\@tempcnta<4095\advance\@tempcnta\@ne\repeat}
% babel-french sets its own classes, and saves and restores the interchar
% state, each time a language is selected: join again after it does.
\addto\extrasfrench{\germain@greekjoin}
\addto\extrasenglish{\germain@greekjoin}
\AtBeginDocument{\germain@greekjoin}
\makeatother

\usepackage{geometry}
\geometry{a4paper,margin=25mm,marginparwidth=28mm}
\usepackage{marginnote}
\usepackage{xcolor}
\usepackage[normalem]{ulem}
\setlength{\emergencystretch}{3em}
\usepackage{hyperref}
\hypersetup{colorlinks=true,linkcolor=black,urlcolor=black,pdfborder={0 0 0}}

% --- Batch metadata ---------------------------------------------------------
% Taken as they stand by the rendering script, which shows them at the head of
% the reading view. They fix what the file claims to cover. The macro names are
% the parent project's, so that the scripts stay shared; \folder here means the
% volume's slug — `fr-9115` — and \pages counts Gallica views.

\newcommand{\folder}[1]{\gdef\grothFolder{#1}}
\newcommand{\batch}[1]{\gdef\grothBatch{#1}}
\newcommand{\pages}[2]{\gdef\grothFirst{#1}\gdef\grothLast{#2}}
\newcommand{\foldertitle}[1]{\gdef\grothFoldertitle{#1}}

% The holder's own shelfmark, printed as they write it: « Français 9115 ».
\newcommand{\shelfmark}[1]{\gdef\grothShelfmark{#1}}

% The dating, copied verbatim from the holder's catalogue — never inferred
% here. « XIXe siècle » is the BnF's; a pass that dates a leaf from its content
% says so in a \note{}, as its own inference.
\newcommand{\dating}[1]{\gdef\grothDating{#1}}

% The status notice, printed in the title block and repeated as a diagonal
% watermark on every page of the PDF. The manuscripts are in the public
% domain; what this line declares is not a right but a quality — an
% unverified first pass by a machine. The skills write the line; a file
% without it gets no watermark, so leaving it out is visible in the output.
\newcommand{\watermark}[1]{\gdef\grothWatermark{#1}}

\gdef\grothFolder{?}\gdef\grothBatch{}\gdef\grothFirst{?}\gdef\grothLast{?}%
\gdef\grothFoldertitle{}\gdef\grothShelfmark{}\gdef\grothDating{}\gdef\grothWatermark{}

% --- The résumé -------------------------------------------------------------
\newenvironment{resume}
  {\small\setlength{\parskip}{0.45em}}
  {\par\medskip\hrule\medskip}

% --- Keywords ---------------------------------------------------------------
% In English, closing the modernised reading's résumé: the modern vocabulary
% under which to search for what the leaves construct. The manifest extracts
% them to tag the volume — this line is the single source of the tags.
\newcommand{\keywords}[1]{\par\smallskip\noindent
  {\footnotesize\textsc{Keywords} --- \textit{#1}}\par}

% --- The critical apparatus -------------------------------------------------

% The Gallica view number, in the margin. This is the anchor that lets a
% disagreement over a formula be settled by looking at *one* image rather than
% a volume of 750. The site uses it to turn the facsimile as the transcription
% is scrolled. A view is one scanned image: usually one side of a leaf.
\newcommand{\page}[1]{%
  \marginnote{\footnotesize\color{gray!70}\textsf{v.\,#1}}%
  \label{page:#1}\ignorespaces}

% The BnF's pencilled foliation, where the leaf carries it and a pass can read
% it: « 198r ». This is what the literature cites — « ff. 198r–208v » — and
% the one line that turns a citation into a view. Recorded, never inferred:
% a folio a pass cannot read is not written.
\newcommand{\folio}[1]{%
  \marginnote{\footnotesize\color{gray!50}\textsf{f.\,#1}}\ignorespaces}

% The modernised reading's coarser counterpart to \page{}: which views a
% section is drawn from.
\newcommand{\pagerange}[2]{%
  \marginnote{\footnotesize\color{gray!70}\textsf{v.\,#1--#2}}%
  \ignorespaces}

% Illegible: never guessed.
\newcommand{\ill}{\textcolor{red!60!black}{[\,\ldots\,]}}

% A doubtful reading: the word is given, and flagged as uncertain.
\newcommand{\uncertain}[1]{\uline{#1}}

% An editorial addition — a missing letter, an implied word.
\newcommand{\add}[1]{\textcolor{blue!50!black}{[#1]}}

% Struck out by the writer. What was crossed out is part of the document.
\newcommand{\struck}[1]{\sout{#1}}

% The transcriber's note, on the state of the page or the mathematics.
\newcommand{\note}[1]{\par\smallskip{\small\color{gray!60!black}\textit{[#1]}}\par\smallskip}

% A marginal note of hers, distinct from the body of the text.
\newcommand{\marginal}[1]{\par\smallskip\noindent\hspace{1em}%
  {\small\color{gray!60!black}#1}\par\smallskip}

% --- Title ------------------------------------------------------------------
% The title block is French, like the documents themselves.

\AddToHook{shipout/background}{%
  \ifx\grothWatermark\empty\else
    \begin{tikzpicture}[remember picture, overlay]
      \node[rotate=52, scale=3.4, align=center, text=gray!18,
            font=\sffamily\bfseries]
        at (current page.center) {\grothWatermark};
    \end{tikzpicture}%
  \fi}

\AtBeginDocument{%
  \begin{center}
    {\large\bfseries Manuscrits de mathématiciens (Gallica) ---
      \ifx\grothShelfmark\empty\grothFolder\else\grothShelfmark\fi}\\[2pt]
    {\itshape\grothFoldertitle}\\[4pt]
    {\small vues \grothFirst--\grothLast
      \ifx\grothBatch\empty\else\ (lot \grothBatch)\fi}\\[2pt]
    \ifx\grothDating\empty\else
      {\small\color{gray!60!black}Datation du catalogue : \grothDating}\\[2pt]
    \fi
    \ifx\grothWatermark\empty\else
      {\small\bfseries\color{red!45!gray}\grothWatermark}\\[2pt]
    \fi
    {\footnotesize\color{gray!60!black}Transcription établie d'après les images IIIF de Gallica.
     Source gallica.bnf.fr / Bibliothèque nationale de France.\\
     Les lectures incertaines sont soulignées, les illisibles notées [\,\ldots\,],
     les ajouts entre crochets ; \textsf{v.} numérote les vues de Gallica,
     \textsf{f.} la foliotation au crayon de la BnF.}
  \end{center}
  \vspace{1em}}

\endinput


\folder{nal-2339}
\shelfmark{NAL 2339}
\batch{3}
\pages{41}{60}
\dating{}
\watermark{Lecture automatique\\non vérifiée}
\foldertitle{Fermat. Mémoires de mathématiques.}

\begin{document}

\page{41}\folio{17r}
\note{Recto. En haut à droite, deux nombres l'un sous l'autre : « 17 » au crayon gris, d'un trait fin, et au-dessous « 17 » au crayon rouge. Voir l'en-tête. Cachet rond « BIBLIOTHÈQUE NATIONALE R.F. » à gauche du titre. Une nouvelle pièce commence en haut de la page, sous un titre ; la main est une cursive posée, à grandes initiales en volutes, qui est celle d'un copiste ou d'un auteur au net : rien sur la page ne permet de le dire.}

\section*{\uncertain{Nouus} Secundarum et Vlterioris ordinis Radicum In Analyticis Vsus}

\uncertain{Nouus} Secundarum et Vlterioris ordinis Radicum In Analyticis Vsus A Domino de Fermat ad Dominum De Carcaui Die 20\textsuperscript{a}. Aprilis Anno 1650 Missus.
\note{Titre en quatre lignes. L'initiale N est une grande volute ; la dernière lettre du premier mot a la forme d'un e plutôt que d'un s (« Nouue ») : « Nouus » est proposé d'après l'accord avec « Vsus » et « Missus ». « Carcaui » touche presque le bord du feuillet ; les dernières lettres se lisent. Le millésime « 1650 » est souligné. C'est la seule date lue dans le lot, et c'est la page, non le catalogue, qui nomme ici Fermat comme auteur et Carcavi comme destinataire.}

Redductio Secundarum et Vlterioris radicum ad primas quæ maximi est In Algebricis momenti Vnicam pro funda-mento Agnosci duplicatæ æqualitatis Analogiam, Eamque quoties opus fuerit Iterandam progressus Ipse quæstionis ostendit.
\note{« Redductio » est écrit avec deux d. La coupure « funda-mento » est celle de la ligne. Les majuscules à l'intérieur des phrases (« In », « Vnicam », « Ipse ») sont celles de la page.}

Proponatur $A$ Cubus $+ E$ Cubo Æquari $Z$ solido.

Item $B$ In $A + E$ q\uncertain{$+$} $D$ in $E$ Æquari $N$ q\textsuperscript{dto}.
\note{après « E q », un petit signe surélevé, qui se lit comme un « + » ou comme le t d'une abréviation (« q\textsuperscript{t} ») ; la reprise plus bas, « $BA + E^2 + DE$ », donne un « + » à cet endroit, et c'est la lecture proposée. « q\textsuperscript{dto} » : quadrato.}

Vt secunda radix deuoluatur ad primam, hæc Sunto præcepta

Quæcunque a secunda Radice adficientur homogenea In Vnam æquationis partem transeunto.

Vt In superiori Exemplo Cum $A^3 + E^3$ Æquetur $Z^{so}$

Ergo $Z^{s} - A^3$ Æquabitur $E^3$.
\note{les exposants sont écrits en chiffres au-dessus de la lettre ; « $Z^{so}$ », « $Z^{s}$ » : Z solido, l'abréviation de la dimension en exposant.}

Similiter Cum $BA + E^2 + DE$ Æquetur $N^2$

Ergo $N^2 - BA$ Æquabitur $E^2 + DE$.

In Vtraque Igitur Æquatione homogenea abs $E$ siue ab secunda Radice adfecta Vnam æquationis partem constituunt.

Si Igitur duplicata huiusmodi æqualitas ad analogiam Reuocetur Erit Vt.

$Z^{s} - A^3$ ad $E^3$, Ita $N^2 - BA$ ad $E^2 + DE$.

Cum Itaque factum sub extremis Comparabitur facto sub mediis tanquam ei Ipsi Æquale, omnia homogenea Diuisionem admittent per $E$ siue per secundam Radicem Vt patet quia secundus et quartus terminus ab $E$ adficiuntur.
\note{les fins de ligne sont remplies de traits. Le texte s'arrête au premier tiers de la page ; le reste est blanc. La suite est au verso, vue 42.}

\page{42}
\note{Verso de la vue 41, sans chiffre en tête (le « 17 » du recto transparaît à l'envers en haut à gauche). Même main ; la pièce continue.}

Erit nempe $Z^sE^2 - A^3E^2 + Z^sDE - A^3DE$ Æle $N^2E^2 - BAE^3$.
\note{« Æle » : un Æ bouclé suivi de « le », abréviation d'« æquale », employée dans toute la pièce. Le « − » devant $A^3DE$ est traversé d'un petit trait vertical, comme un « + » corrigé ou l'inverse ; le calcul demande « − ». « $N^2E^2$ » est écrit tel quel, l'exposant 2 net à côté du 3 de « $BAE^3$ » : le produit des moyens donne $N^2E^3$, et la ligne suivante, divisée par $E$, a bien $N^2E^2$. L'erreur est laissée.}

Omnia diuidantur per $E$ donec aliquod ex homogeneis adfectione sub $E$ continuo Liberetur.

Erit $Z^sE - A^3E + Z^sD - A^3D$ Æle $N^2E^2 - BAE^2$.

Quo peracto noua hæc æquatio Vno ad minus gradu depressior erit (quoad secundam Radicem) quam elatior ex duabus primum propositis

Patet nempe Elatiorem ex duabus primum propositis adfici sub Cubo $E$ Istius Vero nullam ab $E$ adfectionem Excedere $E$ q\textsuperscript{dtum}

Nec tamen sic quiescendum sed Iteranda duplicatæ æqualitis analogia donec adfectio secundæ Radicis fiat tantum sub Latere, Vt Asymmetria omnis Euanescat.
\note{« æqualitis » : ainsi écrit, pour « æqualitatis ».}

Præparetur Itaque vltima hæc Æquatio Iuxta modum præscriptum Vt homogenea sub $E$ quomodocunque affecta Vnam æquationis partem faciant.

Erit Itaque $Z^sD - A^3D$ Æle $N^2E^2 - BAE^2 - Z^sE + A^3E$.

Sed ex duabus primum propositis quæ depressior est exhibet æquationem sequentem.

$N^2 - BA$ Æle $E^2 + DE$.

Reuocetur rursum ad \struck{\ill{}} Analogiam duplicata Illa æqualitas

Erit Itaque $Z^sD - A^3D$. / ad $N^2E^2 - BAE^2 - Z^sE + A^3E$ // vt $N^2 - BA$ / ad $E^2 + DE$.
\note{les termes de la proportion sont séparés par des barres obliques, une simple entre les termes d'un rapport, une double entre les deux rapports ; elles sont gardées.}

Cum Itaque factum sub extremis Æquabitur facto sub mediis \struck{Vt} tanquam Ipsi Æquale, \struck{om} omnia homogenea poterunt diuidi per $E$ Vt supra demonstratum est.

Erit nempe.

$Z^sDE^2 + Z^sD^2E - A^3DE^2 - A^3D^2E$.
\note{la page s'arrête sur ce premier membre, aux deux tiers de sa hauteur ; l'égalité se poursuit en tête de la vue 43.}

\page{43}\folio{18r}
\note{Recto. En haut à droite, « 18 » au crayon gris et au-dessous « 18 » au crayon rouge, comme à la vue 41. Même main ; la page poursuit l'égalité laissée au bas de la vue 42.}

Æquale

$N^4E^2 - N^2BAE^2 - N^2Z^sE + N^2A^3E - BAN^2E^2 + B^2A^2E^2 + BZ^sAE - BA^4E$.

Et omnibus abs $E$ diuisis fiet tandem.

$Z^sDE + Z^sD^2 - A^3DE - A^3D^2$ Æle $N^4E - N^2BAE - N^2Z^s + N^2A^3 - BAN^2E + B^2A^2E + BZ^sA - BA^4$
\note{le second membre est coupé sur deux lignes, la seconde rejetée à droite ; un tiret suivi d'une virgule marque la coupure à gauche.}

Quo peracto Noua hæc Æquatio Vnius adhuc gradus depressionem (quoad secundam Radicem) Lucrata est Vt hic patet Cum enim homogenea sub $E$ adfecta In vnam Æquationis partem transierint fiet.
\note{« peracto » : le a est écrit au-dessus de la ligne. La fin de « transierint » est empâtée d'une tache d'encre.}

$Z^sD^2 - A^3D^2 + N^2Z^s - N^2A^3 - BZ^sA + BA^4$ Æle $N^4E - N^2BAE - BAN^2E + B^2A^2E - Z^sDE + A^3DE$.

Neque vlterius progrediendum Cum Iam secunda radix sub Latere tantum appareat, Ideoque solo Applicationis beneficio Ipsius $E$ relatio ad primam radicem manifestabitur.

Vt fit.
\[ \frac{Z^sD^2 - A^3D^2 + N^2Z^s - N^2A^3 - BZ^sA + BA^4}{N^4 - N^2BA - BAN^2 + B^2A^2 - Z^sD + A^3D} \text{ Æquabit}^{ur}\ E. \]
\note{« Vt fit. » est écrit petit, au-dessus du numérateur, contre la marge. La barre de fraction court sur toute la ligne et « Æquabit\textsuperscript{ur} E » est écrit à son extrémité droite.}

quo tendendum Erat.

Vt Igitur duæ primum propositæ radices In vnam transeant resumatur Ex duabus prioribus Æquationibus quam Volueris depressior tamen Idonea magis ne Alterius asendat Æquatio.
\note{« ne Alterius asendat » : ainsi écrit ; le sens attend « ne vlterius ascendat ».}

Cum Itaque In vna Ex Æquationibus primum propositis $BA + E^2 + DE$ Æquetur $N^2$. Loco Ipsius $E$ subrogetur Iam agnitus eius valor per Relationem Vel ad terminos cognitos vel ad priorem Radicem quæ In Exemplo proposito Est $A$ Et Rursum sub hac noua specie ordinetur æquatio manifestum est euanuisse omnino secundam radicem $E$ In æquatione ab omni Asymmetria Liberam Itum esse methodumque esse generalem.
\note{« Liberam Itum esse » : ainsi écrit. Une tache d'encre couvre le haut de « æquatione ».}

Si enim plures duobus termini proponant\textsuperscript{ur} Incogniti methodus Iterata tertias si opus fuerit Radices ad primas et secundas; Deinde secundas ad primas \&c. eodem prorsus artificio reducet.
\note{la ligne finit par un trait de plume. Le texte s'arrête au milieu de la page, le reste est blanc ; la vue 44 ouvre un « Appendix » à cette méthode.}

\page{44}
\note{Verso de la vue 43, sans chiffre en tête (le « 18 » du recto transparaît à l'envers en haut à gauche). Même main. La page ouvre, sous un titre à grandes volutes, un appendice à la méthode qui précède.}

\section*{Appendix Ad superiorem Methodum}

Superiori methodo debetur perfecta et absoluta asymmetriarum In algebricis Expurgatio; neque enim \uncertain{Symetria} climatismus Vietæa quæ vnicum hactenus ad Asymmetrias \uncertain{huic} Remedium, Efficax satis et sufficiens Inuenta.
\note{« Symetria climatismus Vietæa » : les deux premiers mots se lisent ainsi, lettre à lettre ; l'accord avec « Vietæa » et « Inuenta » attend un nom féminin, que la page ne donne pas clairement. Le dernier mot de la troisième ligne, une majuscule bouclée suivie de « uic », est lu « huic » avec doute.}

Proponatur quippe Latus Cubicum $(BA^2 - A^3) +$ Latus quadratum $(A^2 + ZA) +$ Latus quadrato-quadratum $(B^3A - A^4) +$ Latus quadratum $(\uncertain{G}A - A^2)$ æquari Rectæ $N$.
\note{dans « $BA^2$ » l'exposant est repassé. La lettre qui ouvre la dernière parenthèse, bouclée, se lit G ou D.}

Qua ratione ab asymmetriis huiusmodi Extricabit \uncertain{sese} questionem suam Analysta Vietæus. An non potius dum crescet Labor, crescet difficultas? Et tandem fatigatus et delusus Nouum ab Analytice Lumen Exposcet?
\note{« Quaratione » est écrit en un mot. Le premier mot de la ligne après « Extricabit » se lit « sese » ou « secet ».}

Hoc sane Luculenter Superior methodus subministrat, Vnicum Exemplum Idque breuissimum adiungimus Recluso enim semel fundamento Cætera apertissime manifestantur.

Proponatur Lat. Cub. $(ZA^2) +$ L. Cub. $(A^3 + B^2A)$ æquari $D$.
\note{la première parenthèse porte ici « $ZA^2$ » seul ; la même quantité est écrite « $ZA^2 - A^3$ » quatre lignes plus bas. Transcrit tel qu'écrit.}

Ita primum ordinetur æquatio Vt Vnica Ex Asymmetriis Vnam Illius partem faciat Fiat nempe.

$D -$ Lat. Cub. $(A^3 + B^2A)$ Æle Lat. Cub. $(ZA^2 - A^3)$.
\note{devant « Æle », une première abréviation semblable est noircie d'une tache ou d'une rature.}

Hoc peracto omnes termini Asymmetri a secundis et Vlterioribus si opus fuerit radicibus denominentur Excepto Eo quem Vnicum In vnam Æquationis parte rejecimus.
\note{« In vnam » : le premier jambage de « vnam » est repassé. Le texte s'arrête au milieu de la page ; la suite est au recto suivant, vue 45.}

\page{45}\folio{19r}
\note{Recto. En haut à droite, « 19 » au crayon gris et au-dessous « 19 » au crayon rouge, souligné. Même main ; l'appendice continue.}

Fingatur verbigratia Lat. Cub. $(A^3 + B^2A)$ esse $E$. Hac enim vna ad eam quam Iniungit superior methodus duplicatæ æqualitatis Analogiam deueniemus.

Erit nempe $D - E$ Æ. Lat. Cub. $(ZA^2 - A^3)$.
\note{« Æ. » : le signe d'égalité abrégé, sans le « le » d'« Æle ». Dans « $ZA^2$ », le A est repassé.}

Et omnibus In Cubum ductis

$D^3 + 3DE^2 - 3D^2E - E^3$ Æ $ZA^2 - A^3$.

Sed et ex hypothesi $E^3$ Æ\textsuperscript{tur} $A^3 + B^2A$.
\note{« Æ\textsuperscript{tur} » : æquatur.}

Ergo oritur duplicata æqualitas Vt In Vtraque (Iuxta methodum) termini abs secundâ radice adfecti in Vnam Æquationis partem sunt conijciendi. Erit nempe.

$ZA^2 - A^3 - D^3$ Æ $3DE^2 - 3D^2E - E^3$

Item $A^3 + B^2A$ Æ $E^3$.

Iteretur toties operatio donec secunda Radix ad primam reuocetur, quo peracto, Loco Ipsius $E$ Nouus Ipsius Valor Vsurpetur Vt sub hac Noua specie quæ iis Ex prioribus æqualitatibus ordinetur, omnia constabunt.

Nec mutila adiungo aut moror In superfluis quis Enim non videt singulos terminos Asymmetros posse Eadem raõe, si non sufficiant secundæ radices, tertius, quartus \&c. In Infinitum Insigniri Quo casu quartam siue Vltimam Radicem tanquam secundam considerabis, Reliquas vero tantisper vel pro primis vel pro terminis cognitis habebis donec Vltima Illa omnino \uncertain{Emanuerit} siue ad primas, secundas et Tertias reductæ fuerint simili prorsus artificio tertias ad secundas et primas, ac deinde secundas ad primas Vt Iam sæpius Inculcauimus.
\note{« raõe » : ratione, abrégé par un trait au-dessus du mot. « tertius, quartus » : ainsi écrit. « Emanuerit » se lit ainsi (un m entre E et a) ; le sens attend « euanuerit ».}

Nulla est ergo asymmetria quam non cogat \uncertain{exulerre} hæc methodus, cuius vsus præsertim Eximius Imò.
\note{« exulerre » : lecture lettre à lettre d'un mot que le sens rapproche d'« exulare ». La phrase commencée par « Imò » continue au verso, vue 46. Le texte s'arrête au milieu de la page.}

\page{46}
\note{Verso de la vue 45, sans chiffre en tête (le « 19 » du recto transparaît à l'envers en haut à gauche). Même main ; la phrase de la vue 45 continue.}

Et necessarius Innumerosa potestatum Euanuerint, non deerit vietæum In arithmeticis quæstionibus artificium, et si veris explicari numeris quæstio non possit, proximæ, quantumuis Libuerit, suppetent solutiones, Cum tamen proximas veris solutiones nullo pacto quandiu durauerint Asymmetriæ) consequi possis.
\note{la construction de la phrase, qui reprend « Imò. » à la fin de la vue 45, est donnée telle qu'elle se lit. Une parenthèse fermante suit « Asymmetriæ » ; aucune parenthèse ouvrante n'a été vue.}

Sed et Vlterius Inquirenti obtulit Semira ad Locorum Superficialium plenam et perfectam notitiam Exinde deriuanda methodus, quæ et iis problematis Inseruit In quibus dantur ab Initio plura quam requirat Ipsa problematis construendi determinatio.
\note{« Semira » est écrit en un mot ; le sens le coupe « se mira ».}

Quod Vt Clarius Intelligas, Sunt quædam problemata quæ Vnicam tantum agnoscunt positionem Ignotam, quæ vocari possunt determinata ad differentiam Inter Ipsa et problemata Localia constituendam Sunt alia quædam quæ duas positiones Ignotas habent, et ad Vnicam tantum numquam possunt reduci, Ea problemata sunt Localia, In prioribus Illis Vnicum tantum punctum Inquirimus In Illis Lineam Sed si problema propositum tres Ignotas positiones admittat, problema huiusmodi non Iam punctum dumtaxat aut Lineam tantum sed Integram Superficiem quæstioni Idoneam Inuestigat, Indeque oriuntur Loci ad Superfiæ. \&c. In reliquis.
\note{« Superfiæ. » : ainsi abrégé, pour « superficiem ». Un trait de plume suit « In reliquis ».}

Sicut autem In prioribus data Ipsa sufficiunt ad determinaõne quæstionis, Ita In secundis Vnum datum deest ad determinationem, In tertiis vero duo tantũ data, determinaõne possunt Complere. Ac contrà, potest fieri Vt quemadmodum In his casibus data aut sufficiant aut desint, Ita In plerisque aliis data Ipsa superflua sint et abundent.
\note{« determinaõne » : determinatione(m), abrégé par un trait au-dessus de « na ». « Ac » est écrit en bout de ligne d'une encre plus noire, sur une première lettre repassée. Le texte s'arrête au milieu de la page ; la vue 47 ouvre un « Exemplum ».}

\page{47}\folio{20r}
\note{Recto. En haut à droite, « 20 » au crayon gris et au-dessous « 20 » au crayon rouge. Même main ; l'appendice continue, avec un exemple et une petite figure.}

Exemplo res fiet Euidens.
\note{à droite de ces mots, sur la même ligne, une figure : un segment horizontal lettré A à gauche, B aux deux tiers, C à droite, B et C marqués d'un petit trait vertical.}

In recta $AC$ datâ, datur rectangulũ $ABC$ datur etiam differentia quadratorũ $AB$, et $BC$.

In hoc casu plura patet offerri data quam determinaõ Ideoque Solutio Ipsius quæstionis Exposuit.
\note{« Exposuit » : ainsi écrit.}

Frequentissimus tamen horum problematũ In physicis præsertim et apud artifices est Vsus, eaque omnia per applicaõne simplicem beneficio nostræ methodi Expediuntur; neque recurrendum ad Extrationem Radicum Licet Æquationes ad quasuis potestates ascendant.
\note{un trait de plume suit « ascendant ».}

Proponatur verbi gratia In quadam quæstione

$A^3 + B^2A$ Æ $Z^2D$.

Item etiam (cum ex hypothesi quæstio supponatur esse abundans; has enim quæstiones abundantes sicut Locales \struck{\ill{}} Deficientes appellare consueuimus.)

$G^sA - A^4$ Æ $B^2N$pl.
\note{« abundans », « abundantes » et « Deficientes » sont soulignés. Au-dessus de « Deficientes », un mot bref est biffé. « $G^sA$ » : G solido in A ; « N pl. » : N plano.}

Duplicato hæc Æqualitas ad Analogiam reuocetur, et de præscripta methodo consideretur Vnica nostra radix Ignota quæ In hoc Exemplo est $A$, Sicut in præcedentibus secundam aut Vlterioris ordinis Radicem Considerauimus, et toties Iuxta methodum Iteretur operatio donec adfectio sub $A$ per simplicem aplicaõne possit Expediri siue non tam ad primas radices quam ad Terminos omnino notos reduci patebit solutio problematis simplicissima, nec Analystam deinceps æquationes quadraticæ, Cubicæ quadrato-quadraticæ \&c. remorabuntur.

Lubet Et Coronidis Loco famosi Illius problematis Datis Ellipsi Et puncto Extra Ipsius planum, superficie
\note{la page s'arrête au milieu de sa hauteur, au milieu de la phrase ; la suite est au verso, vue 48.}

\page{48}
\note{Verso de la vue 47, sans chiffre en tête (le « 20 » du recto transparaît à l'envers en haut à gauche). Même main ; la phrase de la vue 47 continue.}

Conicam cuius vertex sit punctũ \uncertain{Datũ} et basis Ellipsis data, Ita plano secare Vt sectio sit Circulus. Solutionem quæ huic methodo debetur, Indicare Eamque simplicissimam.
\note{un trait oblique suit « simplicissimam ».}

Eo deducunt quæstionem Geometræ Vt sumptis quinque punctis ad Libitum In Ellipsi, Vt Iunctis rectis a vertice Conicæ superficiei ad puncta Illa per Iunctas quinque rectas circulum describant.

Sed cum puncta In Ellipsi sint Infinita, si Loco quinque punctorum sumatur sex, fiet problema abundans, et orietur necessario Duplicata Æqualitas quæ tandem Ignotam quantitatem per simplicem Applicationem patefaciet.

Eadem ratione si detur quæcumque Linea curua In plano aut etiam superficies Localis, cuiuscumque tandem gradus sint Inuenientur Diametri et axes figurarum. Imo et In superficie Locali exhibebuntur omnes omnino curuæ Loci superficialis constitutiuæ \&c.

Exponatur verbi gratia superficies Conica cuius vertex sit punctum datum, basis verò parabol\uncertain{e} aut Ellipsis Cubica aut quadrato-quadratica aut Vlterioris in Infinitum gradus.

\ill{} huiusmodi superficies Conica beneficio Illius methodi Ita secari, Vt In Ea Exhibeatur quælibet Curua quæ Ex Constitutione figuræ In ea superficie potest describi, Et problematis solutio semper euadet simplicissima.
\note{les deux mots brefs qui ouvrent la phrase, après une initiale en volute, se lisent lettre à lettre « Oti di » ou « Vti dr » ; la construction (« … secari ») attend un verbe comme « poterit ». Laissés illisibles.}

Nihil addimus de Tangentibus Curuarum Et plerisque aliis huiusmodi Vsibus, fient quippe obuij, nec sedulam Indagatoris Analytici meditationem effugient.
\note{la ligne finit par un trait de plume qui clôt la pièce ; la fin de la page est blanche. Aucune signature. La pièce commencée à la vue 41 (« Nouus Secundarum et Vlterioris ordinis Radicum In Analyticis Vsus ») se termine ici.}

\page{49}\folio{21r}
\note{Recto. En haut à droite, « 21. » au crayon gris et au-dessous « 21 » au crayon rouge. Même main que les vues 41–48, qui passe ici au français. Une lettre commence sous un titre à grande volute.}

\section*{Lettre de Mons.\textsuperscript{r} De fermes A Monsieur De Carcaui}

Lettre de Mons.\textsuperscript{r} De fermes

A Monsieur De Carcaui De Castres

Le 20\textsuperscript{e} auril 1650.
\note{« De fermes » se lit ainsi, en deux mots, la fin nette (« -mes ») ; le titre de la pièce latine, vue 41, porte « de Fermat ». Le « r » de « Mons.\textsuperscript{r} » est pris dans la volute du L. « De Castres » est écrit à la suite du nom du destinataire, séparé de lui par un blanc : la page ne dit pas si c'est le lieu d'où la lettre part ou une qualité du destinataire. La date est celle de la pièce latine (« Die 20\textsuperscript{a}. Aprilis Anno 1650 », vue 41), que la lettre accompagne selon toute apparence (« ma methode generale pour le desbrouilement des Asymmetries ») ; ce rapprochement est une inférence.}

Monsieur

Je n'ay point voulu differer a vous Enuoyer ma methode generale pour le desbrouilement des Asymmetries, Les festes m'ont tout a propos donné le loisir necessaire pour y Vacquer. Je vous enuoye mon original par pure paresse et vous prie me le renuoyer au plutôt, ou bien vne coppie a vostre chois. Vous mesnagerez mes Interestz comme vous l'entendez; Ils consistent seulem.\textsuperscript{t} a me laisser La satisfaction (J'use a dessein de ce mot adouci) d'auoir deuoilé vne matiere qui n'estoit pas cognüe, ce que diuerses questions que Je vo.\textsuperscript{s} ay proposées a diuerses fois, et dont pas vne solution n'a Jamais esté donnée preuuent assés suffisamment.
\note{« tout » est empâté d'une tache. La fin de la phrase est suivie d'un trait oblique.}

Mais si vous voulés auoir le plaisir tout entier proposez hardiment a trouuer La tangente d'vne courbe dont par Exemple La proprieté soit (en prenant $A$ pour L'appliquée, et $E$ pour La portion du diametre qui luy correspond)

Lat. Cub. $(Z^2A - A^3) +$ Lat. quad. quad. $(B\,$pp\textsuperscript{l}$\uncertain{- D^2BA} + A^4)$ $+$ Lat quad. $(BA - A^2) +$ Lat quad. Cub. $(A^5 - B^4A)$.
\note{dans la deuxième parenthèse, « B pp\textsuperscript{l} » (B plano-plano) est souligné ; le terme suivant, une lettre bouclée (D ?) avec un 2 repassé ou barré, puis « BA », est lu avec doute, et le signe qui le précède se confond avec le soulignement.}

hæc omnia quatuor homogenea (quæ in hoc casu sunt Rectæ) æquentur $B + A - E$.

Quæritur tangens ad punctum datum in Curuâ, cuius superior æqualitas proprietatem specificam repræsentat.
\note{la page s'arrête ici, au milieu de sa hauteur ; la lettre continue au verso, vue 50.}

\page{50}
\note{Verso de la vue 49, sans chiffre en tête (le « 21 » du recto transparaît à l'envers en haut à gauche). Même main ; la lettre continue.}

Que fera en ce rencontre La methode de Mons.\textsuperscript{r} Des Cartes que vous sçauës estre Infiniment plus Embarassée que la mienne si les Asymmetries ne sont ostées. pour les oster La methode que Je vous enuoye en vient about sans nulle difficulté. Car en donnant a chacune des Lignes Irrationelles Le nom d'vne seconde Racine. tierce, quarte \&c. on vient tousiours a des doubles Egalitez Lesquelles se reiterent, Jusque a ce que L'application ou la diuision simple oste la derniere de ces racines, puis la penultiesme et ainsy en retrogradant Jusqu'a ce que toutes Les nouuelles Racines Incognues que vous aurés pris a discretion ayent entierem.\textsuperscript{t} disparu, Et pour Lors Il vous restera vne Equation sans Asymmetrie en laquelle Il n'y aura des Racines Incognües que les deux premieres $A$ et $E$, qui n'auront que changé de degré a cause des multiplications frequentes et necessaires a chaque operation, et cette Equation Exempte d'Asymmetries representera la proprieté specifique de la Courbe.
\note{« Mons.\textsuperscript{r} » : le r est porté sur une volute au-dessus de la ligne. « about » est écrit en un mot.}

Or quand nous auons La proprieté specifique de la Courbe sans Asymmetries ma methode de Tangentibus, donne La tangente tres simplement et par La seule application en tous les Cas generalement, soit que La proprieté specifique ait relation a des Lignes droites seulement, soit qu'elle L'ait aussi a des Courbes, et partant en Joignant Les deux methodes, La tangente de la question proposée se trouue par L'application simple, ce qui semble merueilleux.
\note{« nous » : la première lettre, ouverte, se lit n plutôt que v.}

Je n'adiouste pas L'operation Entiere pour ce que La Longueur du trauail me laisseroit, mais en vn mot, Il suffit que vous voyez tres clairement le progrez et La fin de L'ouurage, ce que Je croy auoir esté Incognu
\note{la page s'arrête ici, au milieu de sa hauteur et de la phrase ; la lettre continue au recto suivant, vue 51.}

\page{51}\folio{22r}
\note{Recto. En haut à droite, « 22 » au crayon gris et au-dessous « 22 » au crayon rouge, sur un long trait rouge. Même main ; la lettre continue et s'achève sur cette page. Cachet rond « BIBLIOTHÈQUE NATIONALE R.F. » au bas du texte.}

Jusqu'a present, puisque Monsieur Des Cartes que Je nomme auec tout le respect qui est deub) a la memoire d'vn si grand \struck{homme} et si merueilleux homme, proposoit comme vne difficulté Insurmontable La question suiuante.
\note{une parenthèse fermante suit « deub » ; aucune parenthèse ouvrante n'a été vue. Le premier « homme » est biffé et repassé.}

Estans donnez quatre points Et Courbe En laquelle prenant vn point a discretion Les droites menées de ce point aux quatre donnez fassent vne somme donnée, trouuer vne tangente a quelconque point de cette courbe \&c.

Ainsy que Je puis faire voir par vne de ses Lettres. Pourtant mes methodes Jointes Ensemble en donne la solution simple et L'operaõn se fait en se Joüant.
\note{« operaõn » : operation, abrégé par un trait au-dessus du mot.}

Vous comprenez par la que le principal et plus considerable effect de cette nouuelle methode paroist aux tangentes de toutte sorte de lignes a l'Infiny, puisque les Tangentes s'y trouuent tousiours par applicaõn simple; et apres cela aux questions que J'appelle abondantes, qui se resoluent aussy par la seule diuision et sans aucunes Extractions de Racines \&c.

En voila trop pour vne seconde Lettre, mais Je suis d'humeur a vous faire paroistre ce que peut nostre ancienne Amitié peut estre ces petits Esclaircissemens seruiront a ce qu'il y aura de trop concis dans mon escrit Latin, quoy que Je ne doute point qu'apres que vous et Mess.\textsuperscript{rs} a qui vous Le Communiquerez y auront vn peu resué, Ils n'en trouuent L'Intelligence et La pratique aysée \&c.
\note{« vne seconde Lettre » se lit ainsi. « mon escrit Latin » renvoie, selon toute apparence, à la pièce latine des vues 41–48 ; c'est une inférence. La lettre finit sur « \&c. » au tiers supérieur de la page ; il n'y a ni formule finale ni signature, et le reste de la page est blanc.}

\end{document}
